\documentclass[10pt,letterpaper]{article}
\usepackage[margin=0.65in,headsep=14pt,footskip=18pt]{geometry}
\usepackage[T1]{fontenc}
\usepackage{lmodern,microtype,amsmath,amssymb,booktabs,array,tabularx,enumitem,fancyhdr,xcolor}
\usepackage[colorlinks=true,urlcolor=blue!45!black,linkcolor=blue!45!black]{hyperref}
\hypersetup{pdftitle={CLAS12 reconstruction metrics: a physicist's cheat sheet},pdfauthor={CLAS12 evaluation documentation},pdfsubject={Equations, populations, and interpretation of finding and noise metrics}}
\definecolor{ink}{HTML}{17384A}
\definecolor{muted}{HTML}{475B67}
\definecolor{wash}{HTML}{EDF3F6}
\setlength{\parindent}{0pt}
\setlength{\parskip}{3pt}
\setlength{\emergencystretch}{2em}
\renewcommand{\arraystretch}{1.10}
\setlist[itemize]{nosep,leftmargin=13pt,topsep=2pt}
\pagestyle{fancy}\fancyhf{}
\fancyhead[L]{\small\color{muted}CLAS12 / reconstruction metrics}
\fancyhead[R]{\small\color{muted}Cheat sheet v1 / 10 November 2026}
\fancyfoot[L]{\scriptsize\color{muted}Metric profile: mc\_reconstruction\_v1 \quad Truth policy: primary\_mc\_min4\_v1}
\fancyfoot[R]{\small\thepage\ / 4}
\setlength{\headheight}{13pt}
\newcommand{\metric}[1]{\ifmmode\text{\footnotesize\ttfamily\detokenize{#1}}\else{\footnotesize\ttfamily\detokenize{#1}}\fi}
\newcommand{\field}[2]{\textbf{#1}\newline\metric{#2}}
\newcommand{\heading}[2]{{\Large\bfseries\color{ink}#1}\par{\color{muted}#2}\par\vspace{3pt}}
\newcommand{\labelhead}[1]{\par\vspace{3pt}{\large\bfseries\color{ink}#1}\par}
\newcommand{\note}[1]{\par\noindent\colorbox{wash}{\parbox{\dimexpr\linewidth-2\fboxsep\relax}{#1}}\par}
\newcolumntype{L}[1]{>{\raggedright\arraybackslash}p{#1}}
\newcolumntype{Y}{>{\raggedright\arraybackslash}X}
\begin{document}

\heading{1 / What are we measuring?}{A quick reference for physicists: the equation, its population, and what the number means.}

\textbf{The task.} Given the Central Vertex Tracker (CVT) clusters in an event, track finding groups them into
predicted tracks or leaves them unassigned. We ask how completely and cleanly it
recovers the \emph{eligible Monte Carlo (MC) targets}. We cannot establish the physical validity
of independently reconstructed non-target tracks, so they are neutral to the primary score.

\note{\textbf{Three vocabulary rules.} A ``hit'' here is one \textbf{cluster token},
not a strip or energy deposit. ``Signal'' means \textbf{eligible target support},
not every MC-origin cluster. ``Background'' in legacy labels means
\textbf{non-target}, not necessarily real-data overlay or a fake physical track.}

\textbf{Eligibility.} Keep the existing direct-primary ownership rule (unambiguous,
strict majority of cluster members) and require $\geq4$ owned clusters per MC track
in the \emph{full event}. Secondary/decay tracks and shorter tracks are non-target.
Count before sorting, chunking or padding; retain all clusters and their provenance.

\labelhead{One truth track and one prediction}
Within event $e$, $T_{ei}$ is the set of clusters in eligible truth track $i$;
$P_{ej}$ is the set assigned to prediction $j$. The sets within each collection
are disjoint. Unassigned clusters belong to no $P_{ej}$. Define
\[
 t_{ei}=|T_{ei}|,\qquad q_{ej}=|P_{ej}|,\qquad a_{eij}=|T_{ei}\cap P_{ej}|.
\]
\begin{tabularx}{\linewidth}{@{}L{0.28\linewidth}L{0.22\linewidth}Y@{}}
\toprule Pair quantity & Equation & Read it as\\\midrule
Completeness $c$ & $a/t$ & Fraction of this truth track's clusters recovered by this prediction.\\
Purity $r$ & $a/q$ & Fraction of this prediction belonging to this particular truth track.\\
Intersection-over-union $J$ & $a/(t+q-a)$ & Correct overlap divided by the union of both sets.\\
Dice $D$ & $2a/(t+q)$ & Harmonic mean of pair completeness and purity: $2cr/(c+r)$.\\
\bottomrule
\end{tabularx}
Prediction size $q$ includes \emph{all} assigned clusters: correct, wrong eligible
track, excluded MC and overlay. Zero overlap gives zero for all four quantities.

\labelhead{Two different answers to ``which prediction belongs to this track?''}
\textbf{Dice association $\mathcal M_e$:}
$\displaystyle\arg\max_{\text{one-to-one pairs}}\sum_{(i,j)\in\mathcal M_e}D_{eij}$.
Hungarian matching maximizes total Dice \emph{within each event}; keep only positive
overlaps. There is no quality threshold. Only these predictions are
\emph{MC-associated} in the new hit and contamination metrics.

\textbf{Accepted matching $\mathcal E_e$:}
require $J\geq0.5$, $r\geq0.5$, $c\geq0.5$, then maximize the number of
one-to-one accepted pairs, breaking ties by total IoU. This independent assignment
defines historical track efficiency. It is \emph{not} a cut on the Dice assignment.
At these defaults $J\geq0.5$ already implies both other cuts; all three remain explicit.

\note{\textbf{Three thresholds must not be conflated.} The finding assignment threshold
decides which clusters enter predictions. The match cuts above decide whether a
truth track counts as accepted. The minimum of four defines \emph{truth eligibility};
it is not a minimum predicted-track size or an accepted-match requirement.}

\textbf{Scope/status.} Pages 1--3 describe track finding; page 4 describes binary
noise tagging. The new Dice evaluator is available offline. It does not yet select
training checkpoints: the finding preset still selects signal ARI with a loss
tie-break; noise tagging selects validation loss. A reported metric is not necessarily
the differentiable loss that trains the model.

\textbf{Reference.} This sheet accompanies
\href{mc_reconstruction_v1.pdf}{\texttt{mc\_reconstruction\_v1.pdf}} (exact definitions,
provenance codes, tie conventions and API) and
\href{primary_mc_min4_v1.md}{\texttt{primary\_mc\_min4\_v1.md}} (eligibility).
Executable definitions are in \metric{track_finding_metrics.py} and \metric{noise_metrics.py}.

\newpage
\heading{2 / The six headline quantities}{Read every denominator aloud. ``Efficiency'' and ``purity'' alone are not complete definitions.}

All sums below pool the evaluated dataset. Let
$S_e=\bigcup_iT_{ei}$, $U_e=V_e\setminus S_e$ (non-target),
$R_e=\bigcup_{(i,j)\in\mathcal M_e}P_{ej}$ (associated predictions), and
$B_e\subseteq U_e$ (\emph{verified pure-overlay} clusters). $V_e$ excludes padding.
\[
\begin{aligned}
N_T&=\sum_e\#\{T_{ei}\}, & N_S&=\sum_e|S_e|, & H_A&=\sum_e|R_e|,\\
C&=\sum_e\sum_{(i,j)\in\mathcal M_e}a_{eij}, &
Q_D&=\sum_e\sum_{(i,j)\in\mathcal M_e}D_{eij}, & K&=\sum_e|\mathcal E_e|,\\
H_B&=\sum_e|B_e|, & B_A&=\sum_e|R_e\cap B_e|. &&
\end{aligned}
\]
\begin{tabularx}{\linewidth}{@{}L{0.34\linewidth}L{0.12\linewidth}Y@{}}
\toprule Metric / output field & Equation & Physical question and limit\\\midrule
\field{Mean track Dice (primary)}{mc_track_dice_mean} & $Q_D/N_T$ & How completely and cleanly is the average eligible track reconstructed? Each truth track has one vote; misses contribute zero.\\
\field{Hit completeness}{mc_hit_completeness_global} & $C/N_S$ & What fraction of eligible clusters is recovered in the selected one-to-one associations? Misses and unclaimed split fragments remain in the denominator.\\
\field{Hit purity}{mc_hit_purity_global} & $C/H_A$ & What fraction of clusters in associated predictions belongs to their assigned truth track? Does not measure the fraction of predicted tracks that are genuine.\\
\field{Track efficiency}{track_efficiency_global} & $K/N_T$ & What fraction of eligible truth tracks passes the accepted-match definition? Binary per track; does not require complete recovery.\\
\field{Overlay contamination}{mc_bg_contamination_global} & $B_A/H_A$ & What fraction of associated prediction clusters is verified pure overlay? Only one component of impurity.\\
\field{Overlay leakage to MC}{bg_leakage_to_mc_global} & $B_A/H_B$ & What fraction of all verified pure-overlay clusters enters associated predictions? Same numerator as contamination, different denominator.\\
\bottomrule
\end{tabularx}

\textbf{Interpretation.} Higher is better for the first four; lower for the last
two, interpreted alongside completeness. All lie in $[0,1]$ when defined. Track
efficiency is conditional on \emph{eligible tracks in the evaluated events}; it is
not efficiency relative to all generated particles, geometrical acceptance, PID
purity, or a cross-section correction. Purity can be high while most tracks are missed.

\labelhead{Impurity is broader than overlay contamination}
Let $X=\sum_e|R_e\cap S_e|-C$ count clusters belonging to the \emph{wrong eligible}
track, and $U_A=\sum_e|R_e\cap U_e|$ count associated non-target clusters.
Then $H_A=C+X+U_A$ and
\[
\underbrace{1-C/H_A}_{\metric{mc_hit_impurity_global}}
=\underbrace{X/H_A}_{\metric{mc_other_track_contamination_global}}
+\underbrace{U_A/H_A}_{\metric{mc_nontarget_contamination_global}}.
\]
Excluded MC can increase $U_A$ without increasing $B_A$. Mixed clusters and clusters
with MC present but unresolved composition are counted separately. Neither is
called pure overlay. All clusters have unit weight; no strip-level overlay fraction
is inferred. Even an eligible, mixed-origin cluster counts wholly toward its assigned truth track.

\note{\textbf{Undefined does not mean zero.} No truth tracks: Dice, completeness and
efficiency are null. No Dice associations: purity and associated contamination are
null. No verified overlay: leakage is null. Unknown associated provenance makes
overlay contamination null; unknown provenance anywhere makes leakage null.
The explicit \metric{mc_bg_contamination_lower_bound_global} still gives $B_A/H_A$
when $H_A>0$. A zero lower bound does not establish zero contamination.}

\textbf{Provenance fractions.} For category $g$ (MC-only, overlay-only, mixed,
MC-present/unresolved, unknown), associated fraction is $|R\cap g|/H_A$ and
non-target associated fraction is $|R\cap U\cap g|/H_A$ (dataset-pooled counts).
They appear as \metric{mc_associated_<g>_fraction_global} and
\metric{mc_nontarget_<g>_fraction_global}. Coverage is $(H_A-Z_A)/H_A$ or
$(|V|-Z)/|V|$, with $Z_A,Z$ the corresponding unknown counts; classified unresolved
MC is included in coverage, not claimed to have fully known composition.

\newpage
\heading{3 / Read the plots without importing the wrong meaning}{Counterexamples, historical names, and the averaging convention.}

\begin{tabularx}{\linewidth}{@{}Yrrrr@{}}
\toprule One-event example & Mean Dice & Hit compl. & Hit purity & Track eff.\\\midrule
Two four-cluster truths, both recovered exactly & $1$ & $1$ & $1$ & $1$\\
Same, plus an independent non-target prediction & $1$ & $1$ & $1$ & $1$\\
One four-cluster truth split into two pure pairs & $2/3$ & $1/2$ & $1$ & $1$\\
Two four-cluster truths merged into one prediction & $1/3$ & $1/2$ & $1/2$ & $1/2$\\
One four-cluster truth, all clusters unassigned & $0$ & $0$ & null & $0$\\
One four-cluster truth + four non-target clusters in one prediction & $2/3$ & $1$ & $1/2$ & $1$\\
\bottomrule
\end{tabularx}
\textbf{Why?} A truth can claim only one split fragment; a merged prediction can
serve only one truth. Inclusive match cuts accept the half-track boundary.
Independent non-target predictions can lower historical prediction-level purity.
In the missed-track row,
\emph{legacy signal ARI is still 1}: all truth clusters share the same unassigned label.

\labelhead{Historical dictionary: these names keep their original definitions}
For a single event suppress $e$: $n=\#T$, $m=\#P$, $k=|\mathcal E|$,
$H^+=\sum_j|P_j|$, $U^+=|U\cap\bigcup_jP_j|$, $L=|S\setminus\bigcup_jP_j|$.
Let $F$ count truths overlapping $>1$ prediction and $G$ predictions overlapping
$>1$ truth. \emph{Any} positive overlap counts in $F,G$.
\begin{tabularx}{\linewidth}{@{}L{0.37\linewidth}L{0.19\linewidth}Y@{}}
\toprule Output name & Event equation & What it actually counts\\\midrule
\metric{track_efficiency} / \metric{miss_rate} & $k/n$; $1-k/n$ & Accepted / unaccepted eligible tracks, respectively.\\
\metric{track_purity} / \metric{fake_rate} & $k/m$; $1-k/m$ & Accepted / unmatched predictions; unmatched is not proof of a physical fake.\\
\metric{matched_iou_mean} & $k^{-1}\sum_{\mathcal E}J$ & Mean IoU of accepted pairs only.\\
\metric{matched_purity_mean} & $k^{-1}\sum_{\mathcal E}r$ & Accepted-pair purity; misses and rejected predictions are absent.\\
\metric{matched_efficiency_mean} & $k^{-1}\sum_{\mathcal E}c$ & Accepted-pair \emph{completeness}, not track efficiency.\\
\metric{split_rate} / \metric{merge_rate} & $F/n$; $G/m$ & Overlap multiplicity, before either matching.\\
\metric{background_rejection} & $1-U^+/|U|$ & Non-target clusters left unassigned, including excluded MC.\\
\metric{background_contamination} & $U^+/H^+$ & Non-target fraction of \emph{all} predictions; includes independent non-target tracks.\\
\metric{signal_loss_to_background} & $L/|S|$ & Eligible clusters left unassigned; wrong-track assignments are not counted as lost.\\
\bottomrule
\end{tabularx}

\textbf{ARI: partition agreement, not efficiency.}
For counts $n_{ij}$ on the selected support, row/column totals $n_{i\cdot},n_{\cdot j}$,
and $N=\sum_{ij}n_{ij}$, define
$I=\sum_{ij}\binom{n_{ij}}2$, $A=\sum_i\binom{n_{i\cdot}}2$,
$B=\sum_j\binom{n_{\cdot j}}2$, $Q=\binom{N}{2}$.
\[
\mathrm{ARI}=\frac{I-AB/Q}{(A+B)/2-AB/Q}.
\]
One means identical partitions; zero is the random-partition reference at fixed
group sizes; negative is possible. \metric{ari_signal} uses eligible clusters only:
it ignores non-target contamination. With one truth track, it returns 1 for one
predicted label (even all unassigned), 0 otherwise; no signal gives null.
\metric{ari_with_background} groups all non-target truth into one class, so can
penalize independent non-target tracks. It is null with fewer than two distinct
truth labels. Other degenerate cases follow sklearn's conventions.

\note{\textbf{Averages are part of the definition.}
Global ratio $=\sum_e x_e/\sum_e d_e$; event macro
$=|\mathcal F|^{-1}\sum_{e\in\mathcal F}x_e/d_e$, where $\mathcal F$ contains
defined events. Historical unsuffixed fields above and ARI are event macro;
\metric{_event_macro} makes this explicit. \metric{_global} pools counts;
\metric{matched_<quality>_mean_global} averages all accepted pairs equally.
\metric{mc_track_dice_mean} weights \emph{truth tracks}. One perfect four-cluster
track in one event plus three missed eight-cluster tracks in another gives
Dice $1/4$, event-macro Dice $1/2$, and hit completeness $1/7$. Never average batch means.}

\textbf{Extra diagnostics} (not physical fake identifications; denominator $\sum m$):
\begin{itemize}
\item \metric{unmatched_predicted_fraction_global}: $(\sum m-\sum k)/\sum m$, the global fake-rate meaning.
\item \metric{non_target_only_predicted_fraction_global}: zero-target-overlap prediction count divided by $\sum m$.
\item \metric{mc_overlap_unassociated_predicted_fraction_global}: overlapping predictions not chosen by Dice, divided by $\sum m$.
\end{itemize}

\newpage
\heading{4 / Noise tagging is a different measurement}{A binary decision for each cluster; no track matching is involved.}

\textbf{Labels and decision.} Signal is class 0, noise is class 1. Under the HDF5
policy, signal means a cluster owned by an eligible MC track; non-target clusters
are noise-class examples. The shared full-event cut changes clusters of one-, two-
or three-cluster tracks to noise in both storage paths; legacy ragged data otherwise
retains its existing noise labels. ``Noise'' does not assert an unphysical origin.
For score $s=P_{\rm model}(\mathrm{noise})$, reject if $s>\tau$; ties are retained.
The standard operating point is $\tau=0.5$.

Pool cluster counts across all events, excluding padding. Use semantic subscripts
to avoid the ambiguity of which class is ``positive'':
\[
\begin{array}{c|cc}
 & \text{predicted signal (retained)} & \text{predicted noise (rejected)}\\\hline
\text{true signal} & n_{SS}\;(\mathrm{TN}) & n_{SN}\;(\mathrm{FP})\\
\text{true noise} & n_{NS}\;(\mathrm{FN}) & n_{NN}\;(\mathrm{TP})
\end{array}
\]
The code's TP/TN convention treats \emph{noise} as positive. Let
$N=n_{SS}+n_{SN}+n_{NS}+n_{NN}$.

\begin{tabularx}{\linewidth}{@{}L{0.38\linewidth}L{0.25\linewidth}Y@{}}
\toprule Metric / output name & Equation & Read it as\\\midrule
\field{Signal retention (recall)}{signal_retention / signal_recall} & $n_{SS}/(n_{SS}+n_{SN})$ & Fraction of true signal kept. A cluster-selection efficiency, not finding track efficiency.\\
\field{Signal precision}{signal_precision} & $n_{SS}/(n_{SS}+n_{NS})$ & Signal fraction of the retained sample. A sample purity, without track identity.\\
\field{Noise rejection (recall)}{noise_rejection / noise_recall} & $n_{NN}/(n_{NS}+n_{NN})$ & Fraction of true noise rejected. Not the reciprocal rejection factor.\\
\field{Noise precision}{noise_precision} & $n_{NN}/(n_{SN}+n_{NN})$ & Noise fraction of what was rejected.\\
\metric{accuracy} & $(n_{SS}+n_{NN})/N$ & Fraction of all clusters classified correctly; sensitive to class prevalence.\\
\metric{balanced_accuracy}\newline\metric{macro_recall} & $(R_S+R_N)/2$ & Equal mean of signal/noise recalls. These two names have the same value.\\
\metric{macro_precision} & $(P_S+P_N)/2$ & Equal mean of signal/noise precisions.\\
\metric{all_noise_accuracy} & $(n_{NS}+n_{NN})/N$ & Accuracy of rejecting everything: a trivial baseline.\\
\bottomrule
\end{tabularx}
$R_S,R_N$ denote the two recalls; $P_S,P_N$ the two precisions.
Zero denominators give null; macro quantities require both component values.
All these fractions lie in $[0,1]$. If 95\% of clusters are noise, rejecting
everything gives 95\% accuracy, 0\% signal retention and 50\% balanced accuracy.
Retained-sample purity depends on the signal/noise mixture, even at fixed recalls.

\labelhead{Threshold scans and ranking scores}
\textbf{Signal ROC AUC} (\metric{signal_roc_auc}): with $u=1-s$,
$\Pr(u_S>u_N)+\tfrac12\Pr(u_S=u_N)$ over signal/noise pairs.
It measures ranking: 1 is perfect, 0.5 is random ordering. It is not retention
or purity at $\tau=0.5$. The ROC axes are signal retention and noise leakage
($1-\text{noise rejection}$).

\textbf{Signal average precision} (\metric{signal_average_precision}):
$\mathrm{AP}=\sum_\ell(R_\ell-R_{\ell-1})P_\ell$, along increasing signal recall
as the score cut is relaxed. This is a step-weighted precision--recall summary,
not trapezoidal ROC area; it depends on signal prevalence. Both ranking metrics
are null in this implementation if either truth class is absent.

\note{\textbf{An operating point needs both axes.} Quote signal retention \emph{and}
noise rejection at a specified threshold. The 95\%/99\% retention thresholds and
COATJAVA-matched retention threshold are calibrated on validation and held fixed on
test; achieved test retention can differ. COATJAVA ``noise'' here means its unassigned
clusters, not a dedicated noise classifier. Class-weighted training scores need not
be calibrated physical probabilities. Noise-score quantiles describe score distributions,
not additional reconstruction efficiencies.}

\textbf{Comparisons and uncertainty.} Compare the same truth policy, event cohort,
prediction settings and metric profile.
Report eligible track/cluster counts and provenance coverage; unknown coverage is
not evidence of clean reconstruction. Uncertainty from clustered events or repeated
training seeds is a separate calculation, not supplied by these point estimates.

\end{document}
